\chapter{Modularisation}
\label{chp:modularisation}

This is the first of three chapters that aim to move from writing code
that works to writing software. In this particular chapter you will consider how
to write your code in a structured way.

In this chapter you will cover:
\begin{itemize}
    \item Importing code from python files.
    \item Fragmenting code in to modular components.
\end{itemize}


\section{Tutorial}
\label{sec:modularisation_tutorial}

You will here consider a specific problem of a general type. You will not
concentrate too much on the writing of the code itself. Instead this chapter
concentrates
on how you can write the code as software that will do more than just solve the
specific problem. It will be able to be used for further problems of the same
type.

Consider a Markov chain model of the Board Game “Snakes and Ladders”:
\begin{enumerate}

\item 

what is the shortest number of turns that are possible to win?

\item 

what is the average number of turns?

\end{enumerate}



To solve this problem you will make use of the Python library \texttt{numpy} which is
discussed in the corresponding Chapter~\ref{} to
carry out efficient numerical calculations.


The problem you are considering is in fact an application of a mathematical
object from probability called a Markov Chain which we will not go in to in
detail here however the relevant ideas are that the probability of being in the
100th square after \(k\) turns can be written down as:
\begin{equation*}
\begin{split}
    (\pi P ^ k)_{100}
\end{split}
\end{equation*}

where:
\begin{equation*}
\begin{split}
\pi = (\underbrace{1, 0, \dots, 0}_{100})
\end{split}
\end{equation*}

and \(P\in\mathbb{R}^{100 \times 100}\) where \(P_{ij}\) represents the probability
of being in the \(i\)th square and going to the \(j_th\) square after rolling the
dice.


There are snakes and ladders between the squares as given in
Table~\ref{tbl:snakes_and_ladders}.


\begin{table}[!hbtp]
\begin{tabular}{l|c|c}
\hline

Snake or Ladder
&

From
&

To
\\
\midrule

Ladder
&

3
&

19
\\
Ladder
&

15
&

37
\\

Ladder
&

22
&

42
\\

Ladder
&

25
&

64
\\

Ladder
&

41
&

73
\\

Ladder
&

53
&

74
\\

Ladder
&

63
&

86
\\

Ladder
&

76
&

91
\\

Ladder
&

84
&

98
\\

Snake
&

11
&

7
\\

Snake
&

18
&

13
\\

Snake
&

28
&

12
\\

Snake
&

36
&

34
\\

Snake
&

77
&

16
\\

Snake
&

47
&

26
\\

Snake
&

83
&

39
\\

Snake
&

92
&

75
\\

Snake
&

99
&

70
\end{tabular}
\caption{Snake and Ladders}
\label{tbl:snakes_and_ladders}
\end{table}


The matrix \(P\) will look like:


\[
    P = \left(\begin{array}{@{}*{11}{c}@{}}
0      & 1/6 & 1/6 & 0 & 1/6    & 1/6    & 1/6    & 0      & 0      & \dots  & 0\\
0      & 0   & 1/6 & 0 & 1/6    & 1/6    & 1/6    & 1/6    & 0      & \dots  & 0\\
\vdots & 0   & 0   & 0 & \ddots & \ddots & \ddots & \ddots & \ddots & \ddots & \vdots\\
0      & 0   & 0   & 0 & 0      & 0      & 0      & \dots  & 0      & 0      & 1
\end{array}\right)
\]

Note that because of the ladder on square 3: \(P_{14}=0\) and \(P_{1, 20}=1/6\). The
first row/column of \(P\) corresponds to the state of not being on the board.


A \texttt{csv} file containing this matrix \(P\) can be found at
\texttt{https://zenodo.org/record/4236275}.


To be able to answer the first question you will write a function to compute \(\pi
P ^ k\) for arbitrary \(\pi\), \(k\) and \(P\):




\begin{pyin}
def get_long_run_state(pi, k, P):
    """
    For a Markov chain with transition matrix P and starting state vector pi,
    obtain the state distribution after k steps.

    This is done by computing pi P ^ k

    Parameters
    ----------
    pi : array
        Starting state vector.
    k : int
        Number of iterations.
    P : array
        Transition matrix

    Returns
    -------
    array
        The state vector after k iterations
    """
    return pi @ np.linalg.matrix_power(P, k)
\end{pyin}





For the second question you are going make use of a theoretic result which is
that if \(P\) is of the form:
\begin{equation*}
\begin{split}
P = \begin{pmatrix}
        Q & R \\
        0 & I
    \end{pmatrix}
\end{split}
\end{equation*}

In this case the \textbf{fundamental matrix} is defined by:
\begin{equation*}
\begin{split}
    N = (I - Q) ^ {- 1}
\end{split}
\end{equation*}

The fundamental matrix of an absorbing Markov chains has a number of potential
applications. One of which is to calculate
the expected number of steps for the Markov chain to be absorbed given by:
\begin{equation*}
\begin{split}
    t = N\mathbb{1}
\end{split}
\end{equation*}

where \(\mathbb{1}\) is a column of 1s.


To be able to code this you want to write a function to compute \(t\) but this
requires ``extracting'' \(Q\) from \(P\):




\begin{pyin}
def compute_t(P):
    """
    For an absorbing Markov chain with transition rate matrix this computes the
    vector t which gives the expected number of steps until absorption.

    Note that this does not assume P is in the required format.

    Parameters
    ----------
    P : array
        Transition matrix

    Returns
    -------
    array
        Number of steps until absorption
    """
    indices_without_1_in_diagonal = np.where(P.diagonal() != 1)[0]
    Q = P[indices_without_1_in_diagonal.reshape(-1, 1), indices_without_1_in_diagonal]

    number_of_rows, _ = Q.shape
    N = np.linalg.inv(np.eye(number_of_rows) - Q)
    return N @ np.ones(number_of_rows)
\end{pyin}





You are in fact going to modularise that function. It does 3 things:
\begin{itemize}
\item 

Extracts the matrix \(Q\) from \(P\);

\item 

Computes \(N\);

\item 

Computes \(t\).

\end{itemize}


All of those tasks could be useful in their own right so you are going to break
up that function in to three separate functions:




\begin{pyin}
def extract_Q(P):
    """
    For an absorbing Markov chain with transition rate matrix P this computes the
    matrix Q.

    Note that this does not assume that P is in the required format. It
    identifies the rows and columns that have a 1 in the diagonal and removes
    them.

    Parameters
    ----------
    P : array
        Transition matrix

    Returns
    -------
    array
        The matrix Q
    """
    indices_without_1_in_diagonal = np.where(P.diagonal() != 1)[0]
    Q = P[indices_without_1_in_diagonal.reshape(-1, 1), indices_without_1_in_diagonal]
    return Q


def compute_N(Q):
    """
    For an absorbing Markov chain with transition rate matrix P that gives
    matrix Q this computes the fundamental matrix N.

    Parameters
    ----------
    Q : array
        The matrix Q obtained from P

    Returns
    -------
    array
        The funamental matrix N
    """
    number_of_rows, _ = Q.shape
    N = np.linalg.inv(np.eye(number_of_rows) - Q)
    return N
\end{pyin}





This now allows you to redefine \texttt{compute\_t} in a simpler way:




\begin{pyin}
def compute_t(P):
    """
    For an absorbing Markov chain with transition rate matrix this computes the
    vector t which gives the expected number of steps until absorption.

    Note that this does not assume P is in the required format.
    """
    Q = extract_Q(P)
    N = compute_N(Q)
    number_of_rows, _ = Q.shape
    return N @ np.ones(number_of_rows)
\end{pyin}


\begin{note}
All the code you have written so far is generic in nature so would be better
placed somewhere that it can be used for different project.
\end{note}



You are going to put these three functions (and the necessary \texttt{import numpy as np} statement) in an 
\texttt{absorption.py} file as can be seen in Figure~\ref{fig:absorption.py}.

\begin{figure}[htbp]
\centering
    \includegraphics[width=.7\textwidth]{assets/absorption/main.png}
    \caption{The three modularised function in a python file.}
    \label{fig:absorption.py}
\end{figure}


You will now use everything you have done so far:
\begin{itemize}
\item 

Download, and extract the data available at
\texttt{https://zenodo.org/record/4236275}. Put the
\texttt{main.csv} file in the \texttt{snakes\_and\_ladders}
directory.

\item 

Open a Jupyter notebook in the 
\texttt{snakes\_and\_ladders} directory
and call it \texttt{main.ipynb} also in the

\end{itemize}


This should look like the following:

\begin{raw}
snakes_and_ladders/
    |--- absorption.py
    |--- main.csv
    |--- main.ipynb
\end{raw}


You can now use all of the code we have written in the notebook, first you can
import the functions in \texttt{absorption.py} like any other python library:




\begin{pyin}
import absorption
\end{pyin}





We will also import \texttt{numpy} and use it to read the data file:




\begin{pyin}
import numpy as np

P = np.loadtxt("main.csv", delimiter=",")
\end{pyin}

\begin{note}
The above commands work because the 3 files are all in the same directory.
\end{note}



Now to compute the shortest number of turns:




\begin{pyin}
k = 1
pi = np.zeros(101)
pi[0] = 1
while absorption.get_long_run_state(pi, k, P)[-1] == 0:
    k += 1
k
\end{pyin}





It is possible to arrive at the last square in 6 turns.


Now to compute the average length of the game:





\begin{pyin}
:tags: [nbval-ignore-output]

t = absorption.compute_t(P)
t[0]
\end{pyin}



\begin{raw}
43.49196169497175
\end{raw}




\section{How to}
\label{\detokenize{building-tools/05-modularisation/how/main:how}}\label{\detokenize{building-tools/05-modularisation/how/main::doc}}

\subsection{Import code from python files}
\label{\detokenize{building-tools/05-modularisation/how/main:how-to-import-code-from-python-files}}

Given a \texttt{<file.py>} file in a directory any other python process in the same
directory can import that file as it would a normal library.


\begin{pyin}
import <file>
\end{pyin}



At that stage it is possible to uses any python object (a \texttt{function}, a
\texttt{class}, a
\texttt{variable}) by referring to the \texttt{<file.py>} as library:

\begin{pyin}
<file>.function
<file>.class
<file>.variable
\end{pyin}


\subsection{Break up code in to modular components}
\label{\detokenize{building-tools/05-modularisation/how/main:how-to-break-up-code-in-to-modular-components}}

When modularising code aim to identify specific components of the code that can
be isolated from the rest.
In practice this means writing multiple functions that use the correct inputs
and outputs in chain for an overall goal.


Often this allows you to write a more comprehensive docstring that explains
specific parts of the implemented process.
As an example, consider the problem of wanting to pay a shared bill after
applying a tip, the following function will do this:




\begin{pyin}
def add_tip_and_get_bill_share(total, tip_proportion, number_of_payers):
    """
    This returns the share of a bill to be paid by `number_of_payers`
    ensuring the total paid includes a tip.

    Parameters
    ----------
    total : float
        The total amount of the bill
    tip_proportion : float
        The proportion of the bill that should be added as a tip (a number
        between 0 and 1)
    number_of_payers : int
        The number of people sharing the bill

    Returns
    -------
    float
        The amount each person should contribute
    """
    tip_amount = tip_proportion * total
    total += tip_amount
    return total / number_of_payers
\end{pyin}

You can check that this works:

\begin{pyin}
add_tip_and_get_bill_share(total=100, tip_proportion=0.2, number_of_payers=6)
\end{pyin}





\begin{raw}
20.0
\end{raw}

An improvement of the above would be:

\begin{pyin}
def add_tip(total, tip_proportion):
    """
    This adds the given proportion to a bill total.

    Note that tip_proportion is a number between 0 and 1. A tip_proportion of 0
    corresponds to no tip and a tip_proportion of 1 corresponds to paying the
    total twice.

    Parameters
    ----------
    total : float
        The total amount of the bill
    tip_proportion : float
        The proportion of the bill that should be added as a tip (a number
        between 0 and 1)

    Returns
    -------
    float
        The total value of the bill (including tip)
    """
    tip_amount = tip_proportion * total
    return total + tip_amount


def get_bill_share(total, number_of_payers):
    """
    This returns the share of a bill by dividing the total by the number of
    payers.

    Parameters
    ----------
    total : float
        The total amount of the bill
    number_of_payers : int
        The number of people sharing the bill

    Returns
    -------
    float
        The amount each person should contribute
    """
    return total / number_of_payers
\end{pyin}


Then to use the above you would be able to explicitly write out each step which
ensures that there is clarity in what is being done:

\begin{pyin}
total = add_tip(total=100, tip_proportion=0.2)
get_bill_share(total=total, number_of_payers=6)
\end{pyin}





\begin{raw}
20.0
\end{raw}



\section{Exercises}
\label{\detokenize{building-tools/05-modularisation/exercises/main:exercises}}\label{\detokenize{building-tools/05-modularisation/exercises/main:modularisation-exercises}}\label{\detokenize{building-tools/05-modularisation/exercises/main::doc}}

\begin{enumerate}

\item 

Use the code written in Section~\ref{sec:modularisation_tutorial} to obtain the average time
until absorption from the first state of the Markov chains with the following transition 
matrices:
\begin{enumerate}

\item 

\(P = \begin{pmatrix}1/2 & 1/2 \\ 0 & 1 \end{pmatrix}\)

\item 

\(P = \begin{pmatrix}1/2 & 1/4 & 1/4\\ 1/3 & 1/3 & 1/3  \\0 & 0 & 1 \end{pmatrix}\)

\item 

\(P = \begin{pmatrix}1 & 0 \\ 1/2 & 1/2 \end{pmatrix}\)

\item 

\(P = \begin{pmatrix}1/2 & 1/4 & 1/4\\ 1/3 & 1/3 & 1/3  \\1/5 & 0 & 4/5 \end{pmatrix}\)

\end{enumerate}

\item 

Modularise the following code by creating a function \texttt{flip\_coin} that takes a
\texttt{probability\_of\_selecting\_heads} variable.

\begin{pyin}
import random

def sample_experiment(bag):
    """
    This samples a token from a given bag and then
    selects a coin with a given probability.

    If the sampled token is red then the probability
    of selecting heads is 4/5 otherwise it is 2/5.

    This function returns both the selected token
    and the coin face.
    """
    selected_token = random.choice(bag)

    if selected_token == "Red":
       probability_of_selecting_heads = 4 / 5
    else:
       probability_of_selecting_heads = 2 / 5

    if random.random() < probability_of_selecting_heads:
       coin = "Heads"
    else:
       coin = "Tails"
    return selected_token, coin
\end{pyin}

\item 

Modularise the following code by writing 2 further functions:
\begin{itemize}
\item 

\texttt{get\_potential\_divisors}: A function to return the integers between 2 and
\(\sqrt{N}\) for a given \(N\).

\item 

\texttt{is\_divisor}: A function to check if \(n | N\) (``\(n\) divides \(N\)'') for given
\(n, N\).

\end{itemize}

\begin{pyin}
import math

def is_prime(N):
     """
    Checks if a number N is prime by checking all that positive integers
    numbers less sqrt(N) than it that divide it.

    Note that if N is not a positive integer great than 1 then it does not
    check: it returns False.
    """
    if N <= 1 or type(N) is not int:
        return False
    for potential_divisor in range(2, int(math.sqrt(N)) + 1):
        if (N % potential_divisor) == 0:
            return False
    return True
\end{pyin}


Confirm your work by comparing to the results of using \texttt{sympy.isprime}.

\item 

Write a \texttt{stats.py} file with two functions to implement the calculations of
mean and population variance.


Note that the mean is defined by:
\begin{equation*}
\begin{split}
    \bar x \frac{\sum_{i=1}^{N} x_i}{N}
   \end{split}
\end{equation*}

The population variance is defined by:
\begin{equation*}
\begin{split}
    \sigma ^ 2 = \frac{\sum_{i=1}^{N} (x_i - \bar x) ^ 2}{N}
   \end{split}
\end{equation*}

Use your functions to compute the mean and population variance of the following
collections of numbers:
\begin{enumerate}

\item 

\(S_1=(1, 2, 3, 4, 4, 4, 4, 4)\)

\item 

\(S_1=(1)\)

\item 

\(S_1=(1, 1, 1, 2, 3, 4, 4, 4, 4, 4)\)

\end{enumerate}


Compare your results to the results of using the \texttt{statistics.mean},
and \texttt{statistics.pvariance}. 

\end{enumerate}



\section{Further information}
\label{\detokenize{building-tools/05-modularisation/why/main:further-information}}\label{\detokenize{building-tools/05-modularisation/why/main::doc}}

\subsection{Why modularise?}
\label{\detokenize{building-tools/05-modularisation/why/main:why-modularise}}\label{\detokenize{building-tools/05-modularisation/why/main:id1}}

Best practice when writing code is to break up code in to modular parts. One
guiding principle described in~\cite{fowler2018refactoring}:




\begin{quote}


``Code should be obvious. When someone needs to make a change, they should be
able to find the code to be changed easily and make the change quickly without
introducing any errors.''
\end{quote}






Whilst this guiding principle is ambiguous and all concepts related to clean
code writing and refactoring are not things that can be covered in this book one
specific principle is the one referred to in\ref{martin2009clean}:
\begin{quote}


``Functions should do one thing. They should do it well. They should do it
only.''
\end{quote}



In some texts on code architecture you will see arbitrary rules about how many
lines of code should be in a given function. Having a function with 10 or more
lines of code might indicate that it can be modularised. \textbf{However}, it is
not
recommended to follow such rules as sometimes they might add more complexity than
they remove. Make your code clear and
ensuring your functions do one thing well and one thing only.


\subsection{Why do I get an import error?}
\label{\detokenize{building-tools/05-modularisation/why/main:why-do-i-get-an-import-error}}

The most probable explanation for this is that you are importing a file that is
not in the same directory or that you have not imported the file with the
correct name.


As an example, if your code is in a \texttt{library} directory but that your
notebook is in a \textbf{different} directory then you will get an error as shown
below:





\begin{pyin}
import library
\end{pyin}





\begin{raw}
---------------------------------------------------------------------------
ModuleNotFoundError                       Traceback (most recent call last)
Cell In[1], line 1
----> 1 import library

ModuleNotFoundError: No module named 'library'
\end{raw}






Similarly if you perhaps incorrectly saved your \texttt{library.py} file with a typo in
the name such as: \texttt{librery.py} then you would get the same error.


\subsection{How do I make my file importable from other directories?}
\label{\detokenize{building-tools/05-modularisation/why/main:how-do-i-make-my-file-importable-from-other-directories}}

This falls under the subject matter of ``packaging''. This is not covered in
this book.
